\chapter{The \texttt{WithLibrary.ino} Sketch}

The initial \path{WithLibrary.ino} sketch was a simple duplicate of \path{BareBones.ino}. However, as the LM75 constants and functions are now in the library, we don't want them to be present in the sketch or we will get ``duplicate identifier'' errors from the compiler.

After extracting all the LM75 ``stuff'' to the library header and source files, then deleting it from the \path{BareBones.ino} sketch, we are left with the \func{setup()} and \func{loop()} functions in the main \path{WithLibrary.ino} sketch file.

Listing \ref{lst: WithLibrary.ino - Header} is the header, and can be skipped over!

\begin{lstlisting}[caption={WithLibrary.ino - Header.},label={lst: WithLibrary.ino - Header}, style=myArduino, firstnumber=1]
    //===============================================================
    // LM75 Temperature sensor demo for MakerSpace presentation.
    //
    // Uses the LM75 library, developed for the presentation, to
    // access the LM75 sensor.
    //
    //
    // Norman Dunbar.
    // 16 June 2026.
    //===============================================================
\end{lstlisting}

Listing \ref{lst: WithLibrary.ino - Prologue} is the new prologue. As you can see, there's not a lot left. All we have to declare are the LM75 sensor's seven bit address and the clock speed we want to run it at on the \twi{} bus.

\begin{lstlisting}[caption={WithLibrary.ino - Prologue.},label={lst: WithLibrary.ino - Prologue}, style=myArduino, firstnumber=12]
    #include <LM75_AB.h>

    // Where is my sensor?
    const int LM75_ADDRESS = 0x4F;

    // Bus speed in Hz.
    const uint32_t LM75_CLOCK_HZ = 100000;    // 100 KHz.
\end{lstlisting}

What I hope you will notice, at line 12, is the fact that the previous inclusion of \path{Wire.h} has gone. By following the instruction above to use a library---see Section \ref{section: Using the library} on Page \pageref{section: Using the library}---the correct header files for the library have been added to the sketch.

The \func{setup()} function is next, see Listing \ref{lst: WithLibrary.ino - Setup}. Nothing has changed in this function since the \path{BareBones.ino} sketch. It is identical.


\begin{lstlisting}[caption={WithLibrary.ino - Setup.},label={lst: WithLibrary.ino - Setup}, style=myArduino, firstnumber=24]
    void setup() {
        uint8_t productID;

        // Start up I2C as a controller. Talking to the sensor
        // at 100 KHz.
        LM75_begin(LM75_ADDRESS, LM75_CLOCK_HZ);

        // Need the Serial Monitor for output.
        Serial.begin(9600);
        Serial.println("LM75 Temperature Sensor");
        Serial.print("LM75 Product ID: ");

        productID = LM75_getProductId(LM75_ADDRESS);
        if (productID == LM75_UNKNOWN_ID)
        Serial.println("Unknown");
        else
        Serial.println(productID);
    }
\end{lstlisting}

Finally, the \func{loop()} function, as per Listing \ref{lst: WithLibrary.ino - Loop} and again, nothing has changed.


\begin{lstlisting}[caption={WithLibrary.ino - Loop.},label={lst: WithLibrary.ino - Loop}, style=myArduino, firstnumber=47]
    void loop() {
        // Read the temperature.
        float result = LM75_getCurrentTemperature(LM75_ADDRESS);

        // Print it out with much messing with Serial!
        Serial.print("Temperature: ");
        Serial.print(result);
        Serial.println(" degrees C.");

        // Delay is usually a bad thing, but here we are not doing
        // anything else,
        delay(1500);
    }
\end{lstlisting}

So, that was easy. We started with some fairly generic LM75 accessing functions in the \path{BareBones.ino} sketch, which we tested and found to be working. We extracted those functions to a library and edited the sketch to remove them, changing it only to use the library's header instead of the Wire library header.
